\subsection{Implicaciones para las decisiones administrativas}
Las respuestas distinguen tres cuestiones que deben resolverse juntas: qué población interesa, cómo se selecciona la evidencia y qué incertidumbre permanece. Asegurar presencia de departamentos mediante estratos, concentrar visitas por conglomerados y elegir cada unidad a un intervalo fijo son decisiones de diseño diferentes. La facilidad de reunir participantes por cuotas, juicio, redes de contacto o conveniencia tampoco convierte esas selecciones en muestras probabilísticas \citep{holmesBusiness2023}.

Las fórmulas de tamaño no sustituyen esas decisiones. Elegir un valor crítico de confianza o utilizar una proporción conservadora permite planificar precisión bajo supuestos explícitos; omitir clientes sin conexión o consultar solo a empleados disponibles puede mantener sesgos incluso con muchas observaciones. Por ello, una auditoría de transacciones o una encuesta de satisfacción debe justificar tanto su cantidad de datos como el mecanismo que produjo esos datos \citep{bradley2021}.

De esta relación se desprende que la representatividad es una condición que se examina, no un atributo automático del nombre del diseño. Una decisión defendible declara población, marco, reglas de selección y límites de respuesta, y mantiene separados los resultados muestrales de los parámetros que pretende estimar \citep{wackerly2010}.
